% =============================================================================
% EIT-EMP-2: INTERVENTION EMERGENCE ALGORITHM
% =============================================================================
%
% Erweiterung zu EIT-EMP-1 (Gamma Emergence)
% Spezifiziert den Algorithmus zur Berechnung optimaler Interventionsprofile
%
% Version: 1.0 (Januar 2026)
% =============================================================================

\section{EIT-EMP-2: Intervention Emergence Algorithm}
\label{sec:ie-emp2-algorithm}

\begin{tcolorbox}[colback=purple!5!white, colframe=purple!75!black,
    title=\textbf{Empirical Extension EIT-EMP-2: Intervention Emergence Algorithm}]

\textbf{Extends:} EIT-EMP-1 (Gamma Emergence Specification)

\textbf{Purpose:} EIT-EMP-1 specifies how $\gamma(\Psi)$ is computed.
EIT-EMP-2 specifies how the \textit{optimal intervention vector} $\vec{I}^*$ emerges from $\gamma(\Psi)$.

\textbf{Core Insight:} Interventions are not selected from a menu of types, but \textit{emerge} as optimal profiles in continuous space given the context-dependent complementarity structure.

\textbf{Status:} Empirical extension (algorithmic)

\end{tcolorbox}

%------------------------------------------------------------------------------
\subsection{The Optimization Problem}
\label{subsec:emp2-optimization}
%------------------------------------------------------------------------------

\begin{tcolorbox}[colback=blue!5!white, colframe=blue!75!black,
    title=\textbf{Definition EIT-EMP-2-D1: Intervention Effectiveness Function}]

The total effectiveness of an intervention profile $\vec{I} \in [0,1]^{80}$ in context $\Psi$ is:

\begin{equation}
\boxed{E(\vec{I} \;|\; \Psi) = \underbrace{\sum_{i=1}^{80} w_i \cdot I_i}_{\text{Direct Effects}} + \underbrace{\sum_{i<j} \gamma_{ij}(\Psi) \cdot I_i \cdot I_j}_{\text{Interaction Effects}}}
\label{eq:effectiveness}
\end{equation}

where:
\begin{itemize}[nosep]
    \item $w_i$ = baseline effectiveness weight of dimension $i$ (from literature)
    \item $\gamma_{ij}(\Psi)$ = context-dependent complementarity (from EIT-EMP-1)
    \item $I_i \in [0,1]$ = intensity of intervention on dimension $i$
\end{itemize}

\end{tcolorbox}

\begin{tcolorbox}[colback=blue!5!white, colframe=blue!75!black,
    title=\textbf{Definition EIT-EMP-2-D2: Constrained Optimization Problem}]

The optimal intervention profile $\vec{I}^*$ solves:

\begin{align}
\vec{I}^* &= \arg\max_{\vec{I}} \; E(\vec{I} \;|\; \Psi) \label{eq:objective}\\
\text{subject to:} \quad & I_i \in [0,1] \quad \forall i \in \{1, \ldots, 80\} \label{eq:bounds}\\
& \sum_{i=1}^{80} c_i \cdot I_i \leq B \quad \text{(Budget Constraint)} \label{eq:budget}\\
& \|\vec{I}\|_0 \leq K \quad \text{(Complexity Constraint)} \label{eq:sparsity}
\end{align}

where:
\begin{itemize}[nosep]
    \item $c_i$ = cost of activating dimension $i$
    \item $B$ = total budget available
    \item $K$ = maximum number of active dimensions (sparsity)
\end{itemize}

\end{tcolorbox}

%------------------------------------------------------------------------------
\subsection{Matrix Formulation}
\label{subsec:emp2-matrix}
%------------------------------------------------------------------------------

\begin{tcolorbox}[colback=gray!5!white, colframe=gray!75!black,
    title=\textbf{EIT-EMP-2-D3: Quadratic Form}]

The effectiveness function can be written in matrix form:

\begin{equation}
E(\vec{I} \;|\; \Psi) = \vec{w}^\top \vec{I} + \frac{1}{2} \vec{I}^\top \Gamma(\Psi) \vec{I}
\label{eq:quadratic}
\end{equation}

where $\Gamma(\Psi)$ is the symmetric $80 \times 80$ complementarity matrix:

\begin{equation}
\Gamma(\Psi)_{ij} = \begin{cases}
0 & \text{if } i = j \\
\gamma_{ij}(\Psi) & \text{if } i \neq j
\end{cases}
\end{equation}

\textbf{Properties of $\Gamma(\Psi)$:}
\begin{itemize}[nosep]
    \item Symmetric: $\Gamma_{ij} = \Gamma_{ji}$
    \item Sparse: Most off-diagonal entries are 0 (no interaction)
    \item Context-dependent: Changes with $\Psi$
    \item Sign-constrained: $\Gamma_{FS}(\Psi) < 0$ always (EIT-T9)
\end{itemize}

\end{tcolorbox}

%------------------------------------------------------------------------------
\subsection{The EMERGE Algorithm}
\label{subsec:emp2-emerge}
%------------------------------------------------------------------------------

\begin{tcolorbox}[colback=green!5!white, colframe=green!75!black,
    title=\textbf{Algorithm EIT-EMP-2-A1: EMERGE (Emergent Intervention Selection)}]

\textbf{Input:}
\begin{itemize}[nosep]
    \item Context vector $\Psi \in [0,1]^8$
    \item Budget $B$, Complexity limit $K$
    \item Target outcome (which AWARE components to improve)
\end{itemize}

\textbf{Phase 1: Compute Complementarity Matrix}
\begin{enumerate}[nosep]
    \item Load $\gamma_0$ baseline values (48 parameters)
    \item Load $\beta$ sensitivities (current estimates)
    \item Compute $\gamma_{ij}(\Psi) = \gamma_{0,ij} + \sum_k \beta_{ijk} \cdot \Psi_k$
    \item Construct $\Gamma(\Psi)$ matrix
\end{enumerate}

\textbf{Phase 2: Identify Candidate Dimensions}
\begin{enumerate}[nosep]
    \item Filter dimensions relevant to target outcome
    \item Rank by $w_i$ (baseline effectiveness)
    \item Select top $2K$ candidates
\end{enumerate}

\textbf{Phase 3: Greedy Synergy Selection}
\begin{enumerate}[nosep]
    \item Initialize $\vec{I} = \vec{0}$, Active $= \emptyset$
    \item \textbf{while} $|$Active$| < K$ and budget remains:
    \begin{enumerate}[nosep]
        \item For each candidate $i \notin$ Active:
        \item \quad Compute marginal gain: $\Delta E_i = w_i + \sum_{j \in \text{Active}} \gamma_{ij}(\Psi) \cdot I_j$
        \item \quad Compute marginal cost: $\Delta c_i = c_i$
        \item \quad Compute efficiency: $\eta_i = \Delta E_i / \Delta c_i$
        \item Select $i^* = \arg\max_i \eta_i$
        \item \textbf{if} $\gamma_{i^*j}(\Psi) < -\theta$ for any $j \in$ Active: \textbf{skip} (crowding-out risk)
        \item Add $i^*$ to Active, set $I_{i^*} = 1$
    \end{enumerate}
\end{enumerate}

\textbf{Phase 4: Intensity Optimization}
\begin{enumerate}[nosep]
    \item For active dimensions, solve continuous optimization:
    \item $\vec{I}^*_{\text{Active}} = \arg\max E(\vec{I})$ s.t. budget
    \item Use quadratic programming (QP) solver
\end{enumerate}

\textbf{Output:} Optimal intervention profile $\vec{I}^* \in [0,1]^{80}$

\end{tcolorbox}

%------------------------------------------------------------------------------
\subsection{Crowding-Out Avoidance}
\label{subsec:emp2-crowding}
%------------------------------------------------------------------------------

\begin{tcolorbox}[colback=red!5!white, colframe=red!75!black,
    title=\textbf{Constraint EIT-EMP-2-C1: Crowding-Out Avoidance Rule}]

\textbf{Hard Constraint (from EIT-T9):}

If dimension $i$ involves Financial (F) and dimension $j$ involves Social (S):
\begin{equation}
I_i \cdot I_j > 0 \implies \text{Explicit justification required}
\end{equation}

\textbf{Soft Constraint (Threshold $\theta$):}

For any pair with $\gamma_{ij}(\Psi) < -\theta$ (default $\theta = 0.15$):
\begin{equation}
\min(I_i, I_j) \leq \epsilon \quad \text{(at most one can be strongly active)}
\end{equation}

\textbf{Interpretation:} The algorithm actively avoids combining dimensions that crowd each other out, unless the context $\Psi$ sufficiently weakens the negative interaction.

\end{tcolorbox}

%------------------------------------------------------------------------------
\subsection{Worked Example}
\label{subsec:emp2-example}
%------------------------------------------------------------------------------

\begin{tcolorbox}[colback=yellow!5!white, colframe=yellow!75!black,
    title=\textbf{Example: Retirement Savings Intervention}]

\textbf{Context:}
\begin{itemize}[nosep]
    \item $\Psi = (0.3, 0.7, 0.5, 0.4, 0.6, 0.8, 0.3, 0.5)$
    \item Target: Increase AWARE for long-term financial decisions
    \item Budget: $B = 100$, Complexity: $K = 4$
\end{itemize}

\textbf{Phase 1: Compute $\gamma(\Psi)$}

Selected relevant pairs:
\begin{align*}
\gamma(\text{Self\_KNU\_F\_pos}, \text{Self\_KNU\_F\_neg}) &= -0.50 + 0.05 \cdot 0.7 = -0.465 \\
\gamma(\text{Self\_KNU\_F\_pos}, \text{Comm\_INU\_S\_pos}) &= \frac{1}{4}[-0.30 + 0.25 + (-0.25) + 0] + \beta \cdot \Psi \\
&= -0.075 + 0.02 = -0.055 \\
\gamma(\text{Self\_KNU\_E\_pos}, \text{Self\_KNU\_Ex\_pos}) &= +0.45 + 0.03 \cdot 0.8 = +0.474
\end{align*}

\textbf{Phase 2-3: Selection}

Candidates ranked by $w_i$:
\begin{enumerate}[nosep]
    \item Self\_KNU\_F\_pos ($w = 0.8$) — core target
    \item Self\_KNU\_E\_pos ($w = 0.6$) — emotional resonance
    \item Self\_KNU\_Ex\_pos ($w = 0.5$) — experience/simulation
    \item Comm\_INU\_S\_pos ($w = 0.4$) — social proof
\end{enumerate}

Check crowding-out:
\begin{itemize}[nosep]
    \item (1) + (4): $\gamma = -0.055 > -0.15$ ✓ OK (weak negative)
    \item (2) + (3): $\gamma = +0.474$ ✓ Strong synergy!
\end{itemize}

\textbf{Phase 4: Optimal Intensities}

QP solution: $\vec{I}^* = (0.85, 0.70, 0.65, 0.25, 0, 0, \ldots)$

\textbf{Emergent Intervention Design:}
\begin{itemize}[nosep]
    \item \textbf{Primary:} Future self visualization (Self\_KNU\_F\_pos, $I = 0.85$)
    \item \textbf{Secondary:} Emotional framing of consequences (Self\_KNU\_E\_pos, $I = 0.70$)
    \item \textbf{Tertiary:} Retirement simulation experience (Self\_KNU\_Ex\_pos, $I = 0.65$)
    \item \textbf{Light touch:} Peer comparison (Comm\_INU\_S\_pos, $I = 0.25$)
\end{itemize}

$\Rightarrow$ This profile \textit{emerged} from the context—it was not pre-selected from a typology!

\end{tcolorbox}

%------------------------------------------------------------------------------
\subsection{LLMMC Integration}
\label{subsec:emp2-llmmc}
%------------------------------------------------------------------------------

\begin{tcolorbox}[colback=blue!5!white, colframe=blue!75!black,
    title=\textbf{Protocol EIT-EMP-2-P1: LLMMC for Algorithm Parameters}]

\textbf{What LLMMC Estimates:}
\begin{enumerate}[nosep]
    \item Baseline weights $w_i$ for dimensions without experimental data
    \item Cost estimates $c_i$ for novel intervention mechanisms
    \item Threshold $\theta$ calibration for crowding-out avoidance
\end{enumerate}

\textbf{LLMMC Query Template:}
\begin{quote}
\textit{``Given an intervention targeting [AWARE component] in context [describe $\Psi$]:}
\begin{enumerate}[nosep]
    \item \textit{Rate the standalone effectiveness of dimension $D_i$ (0-1)}
    \item \textit{Estimate implementation cost relative to baseline (0.5-2.0)}
    \item \textit{At what negative $\gamma$ threshold would you avoid combining $D_i$ with $D_j$?''}
\end{enumerate}
\end{quote}

\textbf{Calibration:} LLMMC estimates are validated against known intervention outcomes and updated via Bayesian learning (same as $\beta$ in EIT-EMP-1).

\end{tcolorbox}

%------------------------------------------------------------------------------
\subsection{Computational Complexity}
\label{subsec:emp2-complexity}
%------------------------------------------------------------------------------

\begin{table}[htbp]
\centering
\caption{EMERGE Algorithm Complexity}
\label{tab:complexity}
\renewcommand{\arraystretch}{1.3}
\begin{tabular}{@{}lll@{}}
\toprule
\textbf{Phase} & \textbf{Complexity} & \textbf{Typical Time} \\
\midrule
1. Compute $\Gamma(\Psi)$ & $O(80^2 \cdot 8) = O(51,200)$ & $< 1$ ms \\
2. Filter candidates & $O(80 \cdot \log 80)$ & $< 1$ ms \\
3. Greedy selection & $O(K \cdot 2K) = O(K^2)$ & $< 1$ ms \\
4. QP optimization & $O(K^3)$ & $< 10$ ms \\
\midrule
\textbf{Total} & $O(80^2 + K^3)$ & $< 15$ ms \\
\bottomrule
\end{tabular}
\end{table}

\textbf{Implication:} Real-time intervention design is feasible. The algorithm can be run interactively as context $\Psi$ changes.

%------------------------------------------------------------------------------
\subsection{Integration with EIT Framework}
\label{subsec:emp2-integration}
%------------------------------------------------------------------------------

\begin{tcolorbox}[colback=purple!5!white, colframe=purple!75!black,
    title=\textbf{How EIT-EMP-2 Completes the Emergence Chain}]

\begin{center}
\begin{tikzpicture}[node distance=1.5cm, auto]
    \node[draw, rounded corners, fill=gray!20] (context) {Context $\Psi$};
    \node[draw, rounded corners, fill=blue!20, right=of context] (gamma) {$\gamma(\Psi)$};
    \node[draw, rounded corners, fill=green!20, right=of gamma] (emerge) {EMERGE};
    \node[draw, rounded corners, fill=yellow!20, right=of emerge] (intervention) {$\vec{I}^*$};

    \draw[->, thick] (context) -- node[above] {\tiny EIT-EMP-1} (gamma);
    \draw[->, thick] (gamma) -- node[above] {\tiny EIT-EMP-2} (emerge);
    \draw[->, thick] (emerge) -- (intervention);
\end{tikzpicture}
\end{center}

\textbf{EIT-3 (Cluster Emergence):} The algorithm doesn't use fixed types. However, if many contexts $\Psi$ produce similar $\vec{I}^*$, these form \textit{emergent clusters} in intervention space—recoverable post-hoc, not assumed a priori.

\textbf{EIT-10 (Data-Mode):} When $\beta$ estimates are uncertain, EMERGE can:
\begin{itemize}[nosep]
    \item E-Full: Use QP with estimated $\Gamma(\Psi)$
    \item E-Partial: Use greedy with robust subset of $\gamma_0$
    \item E-None: Use only $w_i$ rankings (no interaction terms)
\end{itemize}

\end{tcolorbox}

%------------------------------------------------------------------------------
\subsection{Summary: The Complete Emergence Pipeline}
\label{subsec:emp2-summary}
%------------------------------------------------------------------------------

\begin{table}[htbp]
\centering
\caption{From Context to Intervention: The Emergence Pipeline}
\label{tab:pipeline}
\renewcommand{\arraystretch}{1.4}
\begin{tabular}{@{}clll@{}}
\toprule
\textbf{Step} & \textbf{Input} & \textbf{Process} & \textbf{Output} \\
\midrule
1 & Situation & Encode as $\Psi \in [0,1]^8$ & Context vector \\
2 & $\Psi$ & $\gamma(\Psi) = \gamma_0 + \sum \beta \cdot \Psi$ & Complementarity matrix \\
3 & $\Gamma(\Psi)$, constraints & EMERGE algorithm & Optimal $\vec{I}^*$ \\
4 & $\vec{I}^*$ & Map to concrete mechanisms & Intervention design \\
5 & Observed outcome & Update $\beta$ via LLMMC & Improved estimates \\
\bottomrule
\end{tabular}
\end{table}

\textbf{Key Properties:}
\begin{itemize}[nosep]
    \item \textbf{No type selection:} $\vec{I}^*$ is computed, not chosen from a list
    \item \textbf{Context-sensitive:} Different $\Psi$ → different $\vec{I}^*$
    \item \textbf{Self-improving:} Each observation refines $\beta$
    \item \textbf{Constraint-respecting:} Crowding-out automatically avoided
\end{itemize}

%------------------------------------------------------------------------------
% END OF EIT-EMP-2
%------------------------------------------------------------------------------
